\clearpage
\section*{Appendix R. Global energy compatibility for compact H8 carriers}
\addcontentsline{toc}{section}{Appendix R. Global energy compatibility for compact H8 carriers}
\noindent\textbf{Scope.} This is a conditional obstruction derived from the
unforced Navier--Stokes energy law. It does not exclude a solution with a
non-negligible exterior correction or a different carrier.

\paragraph{Theorem R.1 (Energy compatibility of global $L^2$ retention).}
Let $u$ solve smooth unforced viscosity-one incompressible Navier--Stokes
on $\mathbb T^3$ for $t<T$. Let $U(t)$ be a nonzero smooth periodic,
divergence-free prepared carrier, $\tau=T-t$, and suppose
\begin{equation}\label{eq:R-energy-profiles}
 \|U\|_2^2=M_0\tau^p(1+o(1)),\qquad
 \|\nabla U\|_2^2=R_0M_0\tau^{p-1}(1+o(1)),
\end{equation}
with $M_0,p,R_0>0$. Assume
\begin{equation}\label{eq:R-laplacian-ratio}
 \sup_{t\in(T-\delta,T)}
 \frac{\|U(t)\|_2\|\Delta U(t)\|_2}
      {\|\nabla U(t)\|_2^2}<\infty
\end{equation}
for some $\delta>0$, and assume the full correction is
\emph{globally} subleading in the energy space:
\begin{equation}\label{eq:R-energy-retention}
 \|u(t)-U(t)\|_2=o(\|U(t)\|_2).
\end{equation}
Then the necessary inequality is
\begin{equation}\label{eq:R-energy-necessary}
 \boxed{R_0\le p/2.}
\end{equation}
In particular, with differentiated similarity asymptotics the
amplitude coefficient
$c_A=\langle U_t-\Delta U+\mathbb P[(U\cdot\nabla)U],U\rangle/
\|U\|_2^2$ satisfies
$\tau c_A\to R_0-p/2$; its positive limit is incompatible
with \eqref{eq:R-energy-retention}.

\paragraph{Proof.}
Put $\varepsilon=\|u-U\|_2/\|U\|_2=o(1)$. Periodic integration
by parts and Cauchy--Schwarz show
\[
 \|\nabla u\|_2\|\nabla U\|_2
 \ge\langle\nabla u,\nabla U\rangle
 =\|\nabla U\|_2^2+\langle u-U,-\Delta U\rangle
 \ge(1-o(1))\|\nabla U\|_2^2,
\]
where \eqref{eq:R-laplacian-ratio} controls the last error.
Thus $\|\nabla u\|_2^2\ge(1-o(1))\|\nabla U\|_2^2$.
Likewise $\|u\|_2^2=M_0\tau^p(1+o(1))$. On the terminal
dyad $I_\tau=[T-\tau,T-\tau/2]$, the exact unforced energy
identity implies
\[
 \begin{aligned}
 M_0(1-2^{-p})\tau^p+o(\tau^p)
 &=\|u(T-\tau)\|_2^2-\|u(T-\tau/2)\|_2^2\\
 &=2\int_{I_\tau}\|\nabla u\|_2^2\,ds\\
 &\ge 2(1-o(1))R_0M_0
  \int_{\tau/2}^{\tau}r^{p-1}\,dr\\
 &=2(1-o(1))\frac{R_0M_0}{p}(1-2^{-p})\tau^p.
 \end{aligned}
\]
Division by the positive common factor gives $1\ge2R_0/p$.
The formula for $c_A$ follows from
$\langle\mathbb P[(U\cdot\nabla)U],U\rangle=0$
and $\langle-\Delta U,U\rangle=\|\nabla U\|_2^2$.
\hfill$\square$

\paragraph{Corollary R.2 (Confinement without numerical quadrature).}
If $U$ vanishes outside an embedded horizontal disk of radius
$R_*\sqrt\tau$, then the scalar disk Dirichlet inequality applied
componentwise yields
\[
 \|\nabla U\|_2^2\ge
 \frac{j_{0,1}^2}{R_*^2\tau}\|U\|_2^2,
\]
where $j_{0,1}$ is the first positive zero of the Bessel
function $J_0$. Therefore if
$j_{0,1}^2/R_*^2>p/2$, the assumptions of Theorem R.1
cannot all hold. For the smooth Appendix O prepared-field
prototype with $h=1/10$, $p=1/5$, $X_*=1$ and
$|\eta|\le1/2$, one has $R_*^2=8/3$ and hence
\[
 R_0\ge (3/8)j_{0,1}^2>2>1/10=p/2.
\]
The contradiction requires no unverified profile quadratures.
The smooth flat compact prototype has a bounded
Laplacian-shape ratio under its differentiated similarity scaling.

\paragraph{Connection with Appendix Q.}
For $w=\alpha U+aW+v$ with $W=\partial_z U$ and
$M\sim M_0\tau^p$, $N=\|W\|_2^2\sim N_0\tau^{-A}$,
$A=\frac12+h$, $D=\frac12-h$, the exact exponent
$A+p=2D$ implies
\[
 \alpha=o(1),\quad a=o(\tau^D),\quad
 \|v\|_2=o(\|U\|_2)
 \quad\Longrightarrow\quad
 \|w\|_2=o(\|U\|_2).
\]
When $v$ is chart-supported in volume $O(\tau^{1+D})$,
the pointwise condition $\|v\|_\infty=o(\tau^{-A})$
implies the displayed global $L^2$ smallness, since
$-A+(1+D)/2=p/2$.
If instead $v$ has a non-negligible \emph{exterior}
$L^2$ component, this theorem does not apply: that
component needs an independent energy, pressure and
complementary-propagator analysis.

\paragraph{Open theorem frontier.}
The theorem rules out globally $L^2$-subleading
unforcing for this confined positive-Rayleigh carrier class,
not an unforced singularity in general. Authentic H8 source
selection, the weighted one-sided seven-ray stress inverse,
a possible exterior energy reservoir, and nonlinear terminal
retention remain unproved. The complete unforced
Navier--Stokes theorem remains OPEN.
